\section{Cote 39 : le théorème de préparation « formel », cas adique et cas discret}\label{sec:39}\verifie{Folder39}

La cote 39, \emph{Fonctions et structures \q{élémentaires} [Schémas formels, théorème de préparation] : notes manuscrites (s.d.)}, datée par l'inventaire \q{[vers 1960-1970]}, compte dix-neuf pages en trois ensembles séparés par des intercalaires vierges : deux rédactions successives d'une axiomatique des classes de séries formelles (pages~2 à~7, puis~9 à~16), et, pages~18 et~19, un théorème énoncé et démontré sous un titre souligné de sa main, \q{Th.\ de préparation “formel”} (\lot{39}{1}{18}). La démonstration commence par une réduction : \q{Comme $A$ est séparé et complet pour sa top.\ linéaire, on est ramené au cas où $A$ est discret}, puis (i) \q{\textbf{signifie} que les $a_i$ ($0 \leq i \leq N-1$) sont nilpotents}. Un premier argument, par Nakayama et le gradué associé, est encadré et barré ; le second considère l'application $\varphi : ((b_i), Q) \mapsto \sum b_i Z^i + QF$ et conclut par le gradué de la filtration $\mathfrak m$-adique, \q{finie}, que $\varphi$ est bijective (\lot{39}{1}{19}).

\begin{enonce}[Théorème de préparation « formel »]
Soient $A$ un anneau linéairement topologisé, séparé et complet, $Z$ une indéterminée, $F = \sum_{i \geqslant 0} a_i Z^i \in A[[Z]]$, et $N \in \mathbb N$ tel que
\begin{enumerate}
\item[(i)] $a_i$ est topologiquement nilpotent pour $0 \leqslant i \leqslant N - 1$ ;
\item[(ii)] $a_N$ est inversible.
\end{enumerate}
Alors $F$ est un élément régulier de $A[[Z]]$, et $A[[Z]]/F A[[Z]]$ est un $A$-module libre de base les images des $Z^i$, $0 \leqslant i \leqslant N-1$ : tout $G \in A[[Z]]$ s'écrit de façon unique $G = QF + \sum_{0 \leqslant i \leqslant N-1} b_i Z^i$, avec $Q \in A[[Z]]$ et $b_i \in A$.
\end{enonce}

\noindent La vérification ne porte pas sur cet énoncé dans sa généralité, mais sur deux cas particuliers : le cas d'un anneau complet pour une topologie $I$-adique, avec $a_i \in I$ pour $i < N$ ; et le cas discret auquel la démonstration du feuillet se ramène, avec $a_0, \dots, a_{N-1}$ nilpotents. Dans les deux cas, on écrit le reste $\sum_{i<N} b_i Z^i$ comme un polynôme $r \in A[Z]$ de degré $< N$.

On s'appuie sur le théorème classique de division de Weierstrass, sous la forme suivante : \emph{soient $A$ un anneau séparé et complet pour la topologie $I$-adique et $F \in A[[Z]]$ ; si l'image $\bar F$ de $F$ dans $(A/I)[[Z]]$ est d'ordre $N$ et si le coefficient $a_N$ de $F$ est inversible dans $A$, alors pour tout $G \in A[[Z]]$ il existe un unique couple $(Q, r)$, $Q \in A[[Z]]$, $r \in A[Z]$ de degré $< N$, tel que $G = QF + r$ ; et si $QF = r$ avec $\deg r < N$, alors $Q = 0$ et $r = 0$.}

\begin{lemme}\label{lem:39-adique}
Soient $I$ un idéal de $A$, $A$ séparé et complet pour la topologie $I$-adique, et $F = \sum a_i Z^i$ avec $a_i \in I$ pour $i < N$ et $a_N$ inversible. Alors $F$ est régulier dans $A[[Z]]$, et tout $G \in A[[Z]]$ s'écrit de façon unique $G = QF + r$ avec $Q \in A[[Z]]$ et $r \in A[Z]$ de degré $< N$.
\end{lemme}

\begin{proof}
Si $I = A$, l'anneau $A$, séparé pour la topologie $A$-adique, est nul, et tout est trivial. Supposons $I \neq A$, de sorte que $A/I$ n'est pas nul. Les coefficients de $\bar F$ d'indice $< N$ sont nuls, puisque $a_i \in I$ ; celui d'indice $N$ est l'image de l'inversible $a_N$, donc est inversible dans l'anneau non nul $A/I$, et en particulier non nul. Ainsi $\bar F$ est d'ordre exactement $N$, et, $a_N$ étant inversible, le théorème de division s'applique. Il donne l'existence et l'unicité de l'écriture $G = QF + r$. Si $FH = 0$, on a $HF = 0$ avec le reste $r = 0$ de degré $< N$, donc $H = 0$ : $F$ est régulier.
\end{proof}

\begin{lemme}\label{lem:39-nilp}
Soit $\mathfrak m$ un idéal nilpotent de $A$. Alors $A$ est séparé et complet pour la topologie $\mathfrak m$-adique. Si $a_0, \dots, a_{N-1}$ sont nilpotents, l'idéal qu'ils engendrent est nilpotent.
\end{lemme}

\begin{proof}
Soit $k$ tel que $\mathfrak m^k = 0$ ; alors $\mathfrak m^n = 0$ pour $n \geqslant k$. Un élément qui appartient à tous les $\mathfrak m^n$ est dans $\mathfrak m^k = 0$ : $A$ est séparé. Si $(x_n)$ est une suite telle que $x_n \equiv x_p \pmod{\mathfrak m^n}$ pour $n \leqslant p$, alors $x_k$ en est une limite : pour $n \leqslant k$, $x_n \equiv x_k \pmod{\mathfrak m^n}$, et pour $n \geqslant k$, $x_n \equiv x_k \pmod{\mathfrak m^k}$, c'est-à-dire $x_n = x_k$. Pour la seconde assertion : l'idéal $\mathfrak m = (a_0, \dots, a_{N-1})$ est contenu dans le nilradical et engendré par un nombre fini d'éléments ; c'est un fait classique qu'un tel idéal est nilpotent (si $a_i^{e} = 0$ pour tout $i$, un produit de $N(e-1)+1$ éléments de $\mathfrak m$ est combinaison de monômes en les $a_i$ dont l'un au moins des exposants atteint $e$). C'est la phrase de la lecture : \q{il est de type fini et engendré par des nilpotents, donc nilpotent}.
\end{proof}

\begin{proof}[Démonstration du cas discret]
Soit $A$ un anneau commutatif quelconque. Supposons que les coefficients $a_i$ d'indice $i < N$ soient nilpotents et que $a_N$ soit inversible. Soit $\mathfrak m$ l'idéal engendré par $a_0, \dots, a_{N-1}$. Par le lemme~\ref{lem:39-nilp}, $\mathfrak m$ est nilpotent et $A$ est séparé et complet pour la topologie $\mathfrak m$-adique. Comme $a_i \in \mathfrak m$ pour $i < N$, le lemme~\ref{lem:39-adique}, appliqué à $I = \mathfrak m$, donne la régularité de $F$ et l'écriture unique.
\end{proof}

\paragraph{Ce que la vérification a établi.} Rien de faux. L'étape discrète est exactement celle de la lecture : un nombre fini de nilpotents engendre un idéal nilpotent, et un anneau est complet pour un idéal nilpotent ; les hypothèses ne servent que comme elles sont écrites, (ii) étant la condition d'inversibilité de la division de Weierstrass, et (i) ce qui rend l'ordre de $F$ modulo l'idéal égal à $N$. Sous la forme adique, aucune finitude de l'idéal n'est requise : la lecture n'emploie \q{de type fini} que pour obtenir $\mathfrak m^j A[[Z]] = \mathfrak m^j[[Z]]$ à l'intérieur de sa propre démonstration par le gradué. Cette démonstration par le gradué associé, qui est celle de la page, n'est pas celle qui a été vérifiée : la vérification passe par la division de Weierstrass, prise comme théorème classique. Ce qui est démontré, en revanche, n'est pas le théorème de la page dans sa généralité. Il l'est dans le cas $I$-adique avec $a_i \in I$, et dans le cas discret ; il ne l'est pas pour un anneau linéairement topologisé quelconque, ni même, dans le cas $I$-adique, sous l'hypothèse exacte (i) : pour la topologie $I$-adique, \q{topologiquement nilpotent} signifie $a_i \in \sqrt I$, condition un peu plus faible que $a_i \in I$, et cet écart n'est pas comblé.

\paragraph{Ce qui reste.} Le cas général d'un anneau linéairement topologisé, séparé et complet, c'est-à-dire le recollement des écritures obtenues sur les quotients discrets $A/J$, $J$ parcourant les idéaux ouverts, qui demande l'identification de $A[[Z]]$ à $\varprojlim_J (A/J)[[Z]]$ ; le cas $I$-adique sous l'hypothèse $a_i \in \sqrt I$. Et tout le reste de la cote : les deux rédactions de l'axiomatique des classes de séries formelles (pages~2 à~16), avec les conditions d'inversion et d'hensélianité, les exponentielles et groupes formels en caractéristique nulle, les sous-schémas formels, la condition sur les coefficients et la condition intégrale. L'issue~\#26 range dans son premier niveau, pour cette cote, \q{Weierstrass division and preparation; henselian pairs; algebraic power series} : la division, sous les formes ci-dessus, est faite ; la préparation au sens d'une factorisation en unité et polynôme distingué, les paires henséliennes et les séries algébriques ne le sont pas.
